% GENERATED FILE. Edit canonical lessons, then run npm run generate:books.
% canonical-source-hash: fnv1a64:8a8188719d74f849
% canonical-lessons: ZH-C77-bukeqi, ZH-C77-mafan, ZH-C77-ganxie, ZH-C77-feichang, ZH-C77-zhen

\chapter{Thanks and Replies}
\label{ch:zh-thanks-and-replies}
\hlchaptermodality{77}

\begin{chapteropening}
\textbf{By the end of this chapter:} \emph{I can answer thanks, thank people warmly and say really and very.}

Complete the last lesson of chapter 77: I can answer thanks, thank people warmly and say really and very.
\end{chapteropening}

\section[\texorpdfstring{\zh{不客气} (bú kèqi)}{不客气 (bú kèqi)}]{\zh{不客气} (bú kèqi) — you're welcome; don't mention it}

\label{lesson:ZH-C77-bukeqi}



\begin{quote}
Before the new one: say the Mandarin for good night, then the Mandarin for bye-bye.
\end{quote}

\subsection*{You'll want to know: \zh{不客气}}
\textbf{\zh{不客气}} — \emph{bú kèqi} — \textquotedblleft{}you're welcome; don't mention it\textquotedblright{}.

\begin{quote}
\textbf{\zh{不客气请进}} \emph{bú kèqi qǐng jìn} — Don't mention it. Do come in!
\end{quote}

bù changes its tone here, rising to bú before the falling kè: bú kèqi.

\begin{cousinweb}
\textbf{\zh{不客气}} (\emph{bú kèqi}) is built from \zh{不} 'not' and kèqi 'polite, formal', from Middle Chinese \emph{khaek} 'guest' and \emph{khj+jH} 'breath' (Guangyun, in Baxter's transcription): no need to be formal.
\end{cousinweb}

\begin{grammarlens}[title={What you've built}]
One more word for thanking and answering thanks.
\end{grammarlens}

\subsection*{Your turn}
\begin{itemize}
\raggedright
  \item \emph{Say it:} \emph{bú kèqi}
  \item \emph{Say it:} \emph{bú kèqi}, once more
  \item \emph{Say it:} \emph{bú kèqi qǐng jìn} — Don't mention it. Do come in!
  \item \emph{Recall:} say the Mandarin for good night, then the Mandarin for bye-bye, then say \emph{bú kèqi} again
\end{itemize}

\begin{tcolorbox}[breakable,colback=teal!4,colframe=teal!35!black,title={Before you move on}]
What does \zh{不客气} mean? (\textquotedblleft{}You're welcome; don't mention it\textquotedblright{}.) Say it once more.
\end{tcolorbox}
\section[máfan]{máfan — to trouble; to bother}

\label{lesson:ZH-C77-mafan}



\begin{quote}
Before the new one: say the Mandarin for bye-bye, then the Mandarin for don't mention it.
\end{quote}

\subsection*{You'll want to know: máfan}
\emph{máfan} — \textquotedblleft{}to trouble; to bother\textquotedblright{}.

\begin{quote}
\emph{bù hǎoyìsi máfan nǐ} — Sorry to trouble you.
\end{quote}

Its characters come later; for now it is a word for the ear.

\begin{cousinweb}
\emph{máfan} is built from má and fán, both inherited from Middle Chinese, \emph{mae} and \emph{bjon} (Guangyun, in Baxter's transcription). The dictionary glosses the first as hemp and the second as 'toil'; how the pair came to mean putting somebody to trouble is not settled.
\end{cousinweb}

\begin{grammarlens}[title={What you've built}]
One more word for thanking and answering thanks.
\end{grammarlens}

\subsection*{Your turn}
\begin{itemize}
\raggedright
  \item \emph{Say it:} \emph{máfan}
  \item \emph{Say it:} \emph{máfan}, once more
  \item \emph{Say it:} \emph{bù hǎoyìsi máfan nǐ} — Sorry to trouble you.
  \item \emph{Recall:} say the Mandarin for bye-bye, then the Mandarin for don't mention it, then say \emph{máfan} again
\end{itemize}

\begin{tcolorbox}[breakable,colback=teal!4,colframe=teal!35!black,title={Before you move on}]
What does máfan mean? (\textquotedblleft{}To trouble; to bother\textquotedblright{}.) Say it once more.
\end{tcolorbox}
\section[gǎnxiè]{gǎnxiè — to thank (warmly)}

\label{lesson:ZH-C77-ganxie}



\begin{quote}
Before the new one: say the Mandarin for don't mention it, then the Mandarin for to trouble somebody.
\end{quote}

\subsection*{You'll want to know: gǎnxiè}
\emph{gǎnxiè} — \textquotedblleft{}to thank (warmly)\textquotedblright{}.

\begin{quote}
\emph{gǎnxiè dàjiā} — Thank you, everyone.
\end{quote}

Its characters come later; for now it is a word for the ear.

\begin{cousinweb}
\emph{gǎnxiè} is built from gǎn 'to be moved' and \zh{谢} 'to thank', both inherited from Middle Chinese, \emph{komX} and \emph{zjaeH} (Guangyun, in Baxter's transcription).
\end{cousinweb}

\begin{grammarlens}[title={What you've built}]
One more word for thanking and answering thanks.
\end{grammarlens}

\subsection*{Your turn}
\begin{itemize}
\raggedright
  \item \emph{Say it:} \emph{gǎnxiè}
  \item \emph{Say it:} \emph{gǎnxiè}, once more
  \item \emph{Say it:} \emph{gǎnxiè dàjiā} — Thank you, everyone.
  \item \emph{Recall:} say the Mandarin for don't mention it, then the Mandarin for to trouble somebody, then say \emph{gǎnxiè} again
\end{itemize}

\begin{tcolorbox}[breakable,colback=teal!4,colframe=teal!35!black,title={Before you move on}]
What does gǎnxiè mean? (\textquotedblleft{}To thank (warmly)\textquotedblright{}.) Say it once more.
\end{tcolorbox}
\section[fēicháng]{fēicháng — extremely; very}

\label{lesson:ZH-C77-feichang}



\begin{quote}
Before the new one: say the Mandarin for to trouble somebody, then the Mandarin for to thank.
\end{quote}

\subsection*{You'll want to know: fēicháng}
\emph{fēicháng} — \textquotedblleft{}extremely; very\textquotedblright{}.

\begin{quote}
\emph{fēicháng gǎnxiè} — Thank you so much.
\end{quote}

Its characters come later; for now it is a word for the ear.

\begin{cousinweb}
\emph{fēicháng} is built from fēi 'not' and cháng, both inherited from Middle Chinese, \emph{pj+j} and \emph{dzyang} (Guangyun, in Baxter's transcription). The dictionary glosses the first as 'it is not so'; the second as 'usual' is not given there, so 'out of the ordinary' is likely but not settled.
\end{cousinweb}

\begin{grammarlens}[title={What you've built}]
One more word for thanking and answering thanks.
\end{grammarlens}

\subsection*{Your turn}
\begin{itemize}
\raggedright
  \item \emph{Say it:} \emph{fēicháng}
  \item \emph{Say it:} \emph{fēicháng}, once more
  \item \emph{Say it:} \emph{fēicháng gǎnxiè} — Thank you so much.
  \item \emph{Recall:} say the Mandarin for to trouble somebody, then the Mandarin for to thank, then say \emph{fēicháng} again
\end{itemize}

\begin{tcolorbox}[breakable,colback=teal!4,colframe=teal!35!black,title={Before you move on}]
What does fēicháng mean? (\textquotedblleft{}Extremely; very\textquotedblright{}.) Say it once more.
\end{tcolorbox}
\section[zhēn]{zhēn — really; truly}

\label{lesson:ZH-C77-zhen}



\begin{quote}
Before the new one: say the Mandarin for to thank, then the Mandarin for extremely.
\end{quote}

\subsection*{You'll want to know: zhēn}
\emph{zhēn} — \textquotedblleft{}really; truly\textquotedblright{}.

\begin{quote}
\emph{nǐ zhēn hǎo} — You're so kind.
\end{quote}

Its characters come later; for now it is a word for the ear.

\begin{cousinweb}
\emph{zhēn} is inherited from Middle Chinese \emph{tsyin} (Guangyun, in Baxter's transcription), glossed 'true, as against false'.
\end{cousinweb}

\begin{grammarlens}[title={What you've built}]
One more word for thanking and answering thanks.
\end{grammarlens}

\subsection*{Your turn}
\begin{itemize}
\raggedright
  \item \emph{Say it:} \emph{zhēn}
  \item \emph{Say it:} \emph{zhēn}, once more
  \item \emph{Say it:} \emph{nǐ zhēn hǎo} — You're so kind.
  \item \emph{Recall:} say the Mandarin for to thank, then the Mandarin for extremely, then say \emph{zhēn} again
\end{itemize}

\begin{tcolorbox}[breakable,colback=teal!4,colframe=teal!35!black,title={Before you move on}]
What does zhēn mean? (\textquotedblleft{}Really; truly\textquotedblright{}.) Say it once more.
\end{tcolorbox}
